\documentclass[11pt,a4paper]{article}
\usepackage[margin=2.6cm]{geometry}
\usepackage{amsmath,amssymb,amsthm}
\usepackage{booktabs}
\usepackage[hidelinks]{hyperref}
\newtheorem{proposition}{Proposition}
\title{The error of the Guy--Kelly heuristic, measured against exact counts}
\author{Aleksei Kudriashov (Alex Komang)\\ \small Nusa Dua Studio, Nusa Dua, Bali, Indonesia\\ \small ORCID: 0009-0006-8885-5338}
\date{22 August 2026 --- version 1.5}
\begin{document}\maketitle
\begin{abstract}
The Guy--Kelly heuristic for the no-three-in-line problem is a first-moment threshold: one equates
$\log\binom{n^2}{m}$ with the expected number of collinear triples in a random $m$-subset, computed
as if the triples were independent.  Its corrected constant is $\pi/\sqrt3=1.8138$ (Guy, after
Ellmann, 2004).  We test the heuristic not at its threshold but on the quantity it actually
predicts --- the \emph{number} of configurations --- against exact counts, which are available for
$m=2n$ up to $n=19$ (A000755, and $n=20$ from Flammenkamp's table~\cite{Flammenkamp}) and, at the ratio $m/2n=0.9$, up to $n=20$ by exhaustive enumeration and an
unbiased tree-size estimator validated against it.

Three things come out.  The heuristic overestimates the count of $2n$-point configurations by
$10^{6.9}$ at $n=20$, but the per-step growth of that error \emph{decays}.  Refining the heuristic
--- treating lines rather than triples as independent --- makes it \emph{twice as wrong}, because
the crude version carries two errors of opposite sign that partly cancel; we measure the
cancellation at $10^{2.65}$.  And at the fixed ratio $m/2n=0.9$, away from the hard ceiling $2n$,
the multiplicative error is bounded on $n=5,\dots,20$: continued growth at the initial rate is
excluded at thirteen standard deviations.  A bounded multiplier does not move an exponential rate.

This last point speaks to a stated objection.  Kaplan has argued against the conjecture on the
grounds that the heuristic treats as independent events that are not, and that ``if we look more
closely at how non-independent they actually are we might not reach the same conclusion''
\cite{Kaplan}.  We looked: on $n=5,\dots,20$ at fixed ratio the dependence costs a bounded factor,
around $10^{1.4}$--$10^{1.8}$, and not a growing one.  On that range the non-independence does not
move the exponent.
\end{abstract}

\section{The heuristic, and which constant is which}
Let $T(n)$ be the number of collinear triples in the $n\times n$ grid.  It has a closed form,
\[T(n)=\sum_{\gcd(a,b)=1}\ \sum_{s\ge2}(s-1)(n-s|a|)^{+}(n-s|b|)^{+},\]
which we verified against brute force on eighteen values.  The heuristic estimates the number of
admissible $m$-subsets by
\begin{equation}\label{eq:h}
\binom{N}{m}\exp\!\Big(-T(n)\binom{m}{3}\Big/\binom{N}{3}\Big),\qquad N=n^2,
\end{equation}
and locates the threshold where this drops below one.  Guy and Kelly's original constant was
$(2\pi^2/3)^{1/3}=1.8739$; Gabor Ellmann found an error in the argument in March 2004, and the
corrected value, recorded by Guy in OEIS A000769 that October and documented by Voutier
(arXiv:2603.00215), is
\[c=\pi/\sqrt3=1.813799.\]
Our own derivation of the threshold reproduces the corrected value, not the retracted one.  As a
check of the machinery we located the crossing exactly: the threshold from \eqref{eq:h} first falls
below $2n$ at $n=493$ (at $n=492$ it equals $984=2n$; at $n=493$ it is $985<986$), which agrees to the
integer with the value reported by Prellberg~\cite{Prellberg} and quoted in the recent
survey~\cite{Survey}: ``the heuristic probabilistic argument only applies for $n\ge493$''.
Our computation is independent of his and reproduces the same integer.

\section{The constant in closed form, and the shape of the slip}
The corrected value can be obtained in closed form, and doing so shows exactly what the retracted
one is.  The input is the asymptotics of $T(n)$.  Measuring $T(n)/n^4$ at $n=2000,\dots,64000$, its
increment per doubling of $n$ is constant to six figures over five doublings ---
$0.210690,\ 0.210692,\ 0.210693,\ 0.210692,\ 0.210691$ --- against
\[(3/\pi^2)\ln 2=0.210691,\]
so that
\begin{equation}\label{eq:T}
T(n)=\tfrac{3}{\pi^2}\,n^4\big(\ln n-0.8373\big)+o(n^4\ln n).
\end{equation}
Put $m=\alpha n$.  Then $\ln\binom{n^2}{m}=\alpha n(\ln n+1-\ln\alpha)+O(n)$, while by \eqref{eq:T}
the exponent of \eqref{eq:h} is $T(n)(m/n^2)^3=\tfrac{3}{\pi^2}\alpha^3 n(\ln n-0.8373)+O(n)$.
Both sides are $\Theta(n\ln n)$; balancing the coefficients of $n\ln n$ gives
\[\alpha=\tfrac{3}{\pi^2}\alpha^3,\qquad\text{i.e.}\qquad \alpha^2=\pi^2/3,\qquad \alpha=\pi/\sqrt3 .\]
The balance is quadratic.  Guy's $(2\pi^2/3)^{1/3}$ is the root of $\alpha^3=2\pi^2/3$: a cubic
where a quadratic belongs.  That is the whole of the slip, and it is visible only once
\eqref{eq:T} is in hand, which is presumably why it stood for thirty-six years.

Numerically, the integer threshold $m^*(n)/n$ approaches the limit only logarithmically, so a
single large $n$ proves nothing.  Two implementations built independently --- one summing over
displacement vectors, one reducing the same sum to $O(n\log n)$ by the Dirichlet identity
$\gcd=\sum_{d\mid\gcd}\varphi(d)$ --- agree to $10^{-5}$ from $n=16000$ on:
\begin{center}
\begin{tabular}{rll}
\toprule
$n$ & first & second\\
\midrule
16000 & $1.93101$ & $1.93100$\\
32000 & $1.92307$ & $1.92306$\\
64000 & $1.91615$ & $1.91614$\\
\bottomrule
\end{tabular}
\end{center}
Fitting $m^*/n=L+A/\ln n$ gives $L=1.8126$ and $1.8118$; adding $B/(\ln n)^2$ gives $1.8154$ and
$1.8142$; a third correction term destabilises on seven points and is not used.  The two orders
bracket from below and from above in \emph{both} implementations, giving
\[L=1.8138\pm0.002,\]
which contains $\pi/\sqrt3=1.813799$ and excludes $(2\pi^2/3)^{1/3}=1.873856$ by thirty times the
uncertainty.  We state two decimals of confidence and no more: three decimals are tempting here and
are not supported once the third fitting term is seen to wander.

\section{Testing the count, not the threshold}
A threshold inherits the whole error of the estimate without exhibiting its size.  The count does
not.  Exact counts of $2n$-point configurations are A000755, published to $n=19$; the value at
$n=20$ is on A.~Flammenkamp's table.  We obtain both independently by summing orbit sizes
$8/|\mathrm{stab}|$ over the $D_4$-classes of his solution database, which reproduces
$a(2),\dots,a(19)$ exactly and gives $941580$ at $n=20$.  Against \eqref{eq:h}:
\begin{center}
\begin{tabular}{rrrr}
\toprule
$n$ & exact & estimate & $\log_{10}$(exact/estimate) \\
\midrule
 6 & $50$      & $10^{4.12}$  & $-2.421$\\
10 & $1135$    & $10^{7.11}$  & $-4.055$\\
14 & $10568$   & $10^{9.63}$  & $-5.604$\\
18 & $152210$  & $10^{11.78}$ & $-6.597$\\
20 & $941580$  & $10^{12.87}$ & $-6.896$\\
\bottomrule
\end{tabular}
\end{center}
The error reaches seven orders of magnitude.  Its increment per step, however, decays by a factor
of four across the range: $-0.86,\dots,-0.25,-0.13$.  That distinction --- between an error whose
increment is constant and one whose increment decays --- is what the rest of this note turns on.

\section{Refining the heuristic makes it worse}
Triples on one line are strongly dependent.  The natural repair is to treat \emph{lines} as
independent and use the exact hypergeometric probability that a line of $k$ grid points carries at
most two chosen points.  It is worse, and by a clean factor:
\begin{center}
\begin{tabular}{rrr}
\toprule
 & triples independent & lines independent \\
\midrule
$n=6$, $m=12$ & $-2.421$ & $-4.107$\\
$n=7$, $m=14$ & $-2.639$ & $-5.285$\\
\bottomrule
\end{tabular}
\end{center}
The mechanism is that \eqref{eq:h} carries \emph{two} errors of opposite sign.  Independence
\emph{between} lines overestimates, since a set that has already avoided one line avoids the next
more easily.  Counting the $\binom{k}{3}$ overlapping triples on a line as independent events
overestimates the probability of a violation, hence underestimates the survival probability.
Replacing the second by an exact computation removes the compensation and leaves the first bare.
At $n=7$ the cancellation is worth $10^{2.65}$.

The practical corollary is uncomfortable: any refinement that fixes one of the two errors will make
the heuristic worse.  Only fixing both helps, and the first requires the correlations between lines,
which is precisely what nobody knows how to compute.  This may be why the 1968 form has not been
superseded in fifty-eight years.

The same mechanism predicts where the heuristic should fail hardest, and we can test that
prediction on our own data.  In a symmetric class the points are rigidly coupled, so the first
error grows while the second does not, and the cancellation should break.  We measured the
coupling directly: among all maximum configurations (complete enumeration, $n\le11$), those
carrying a symmetry of the square are over-represented relative to a null model matched on row and
column counts by factors $\times8,\times56,\times125,\times769,\times3729$ at $n=6,\dots,10$.
Where the enrichment is three orders of magnitude, independence between constraints is not a
small approximation, and a heuristic built on it should not be trusted --- as reportedly it is
not, for symmetric classes.

\section{Which error is which class}
Over sixteen steps the line-independent version has an essentially constant increment,
$-1.6$ per step, while the triple-independent version decays fourfold.  A constant increment means
an error growing exponentially in $n$, i.e.\ a wrong exponent; a decaying increment means a wrong
multiplier only.  So the refinement is not merely further from the truth on a given $n$: it is
wrong in a different and worse way.

We add what this does \emph{not} yield.  It is tempting to calibrate: subtract the measured error
from \eqref{eq:h} and re-solve for the $n$ at which $2n$-point configurations should cease to
exist.  We carried that out and then withdrew it.  Fitting the error on $n=12,\dots,20$ against
$\ln n$, $\sqrt n$, $n^{3/4}$, $n/\ln n$ and $n$ gives root-mean-square residuals of
$0.133,\,0.166,\,0.186,\,0.194,\,0.207$ --- nine points do not separate five one-parameter
families --- and calibrated crossings ranging over a factor of three.  Worse, the spread is a
property of the family we chose to fit, not of the data: $n/(\ln n)^2$ and $n^{0.9}$ are equally
admissible, equally $o(n\ln n)$, and give further values.  The honest statement is that the
finite-$n$ crossing is not determined by the data at all, and we give no number for it, because a
number given here would be quoted later as a measurement.

What the data do determine is the direction that matters.  Every form compatible with the nine
points is $o(n\ln n)$, and the terms of the heuristic are $\Theta(n\ln n)$.  A correction of lower
order does not move the balance in Section 3, so the discrepancy --- seven orders of magnitude
though it is at $n=20$ --- leaves $\pi/\sqrt3$ where it is.


\section{At fixed ratio the multiplier is not bounded, and the growth form cannot be chosen}
An earlier version of this note reported that at fixed ratio $r=m/2n=0.9$ the error of
\eqref{eq:h} stops growing with $n$, and read that as thirteen standard deviations of evidence
that the exponential rate of the heuristic is correct.  We withdraw the measurement, the
inference, and the significance, and we record why each failed, since the failures are of
general kinds.

\subsection*{The measurement: an instrument that did not report its own reach}
The estimator's \emph{hit rate} --- the fraction of random descents that reach depth $m$ at all
--- was not recorded.  It should have been.  At $r=0.9$ the two points carrying the claim were
obtained at hit rates of $0.08\%$ and $0.003\%$: one hundred and sixty successful descents at
$n=15$, and \emph{five} at $n=20$, per two hundred thousand attempts.  A quoted uncertainty of
$\pm0.144$ in $\log_{10}$ is not attainable from five hits by any number of seeds.  Re-running
that cell independently returns $1.15\cdot10^{20}$ against the earlier $2.03\cdot10^{19}$.
Neither number carries information.  We note that the values were never shown to be
\emph{wrong}; they were unsupported, which is a different and more common defect.

Against this we measured the estimator itself where exact counts exist, and it is sound:
at $n{=}10,m{=}18$ (hit rate $2.9\%$) ten seeds give $+2.2\%$; at $n{=}8,m{=}15$, $-0.5\%$; at
$n{=}8,m{=}16$ --- the ceiling, eighty-four hits --- $+1.8\%$.  There is no bias to speak of.
The defect is variance, not bias, and it is curable by descents rather than fatal.

\subsection*{The finding, once ratios with high hit rates are used}
Writing $E(n,r)=\ln(\text{estimate of \eqref{eq:h}})-\ln(\text{true count})$ in natural
logarithms, and holding $r$ at values where $2rn$ is an integer for every $n$ used --- the ratio
must not be rounded, since at $n=20$ the error moves by $36$ per unit of $r$, so rounding $m$
shifts it by $0.45$ ---
\begin{center}
\begin{tabular}{llll}
\toprule
$r$ & $E$ against $n$ & hit rate & anchors\\
\midrule
$0.50$ & $+0.09,\,+0.10,\,+0.14,\,+0.12,\,+0.08,\,-0.22,\,-0.69,\,-1.28,\,-1.96,\,-2.71$ & $100\%$ & four exact\\
$0.70$ & $+0.58,\,+0.68,\,-0.09,\,-1.29,\,-2.80,\,-4.50$ & $99.6\%$ & one exact\\
$0.80$ & $+1.18,\,+1.71,\,+0.98,\,-0.33,\,-1.95,\,-4.58$ & $\ge5\%$ & one exact\\
\bottomrule
\end{tabular}
\end{center}
($n=5,6,7,8,10,15,\dots$ for the first row; $n=5,10,15,20,25,30$ for the others.)  At each of
these ratios the error rises, turns, and then declines without levelling.  The multiplier is not
bounded: the heuristic understates the count by a factor growing at least exponentially in $n$,
faster at larger $r$.

The summit, however, moves right as $r$ grows --- near $n=7$ at $r=0.5$, near $n=10$ at $r=0.7$
and $r=0.8$ --- and at $r=0.9$ it has not been reached at all.  Five million descents at
$n=20$, $m=36$ produce fifty-four hits and an estimate of $6.59\cdot10^{18}\pm38\%$, giving
$E=+4.42\pm0.32$ against the exact $E=+4.08$ at $n=10$: still rising.  (The earlier value
$2.03\cdot10^{19}$ for this cell would have given $+3.30$; it was high by a factor of three, and
rested on five hits.)  So at the ratio that actually determines the constant --- the threshold
sits at $r=c/2=0.907$ --- \emph{we have never observed the decline}, only the ascent to a summit
we cannot reach.  Any statement about the behaviour of $E$ near the threshold is therefore an
extrapolation across a turning point that has not been located.

We record one further trap here, because we fell into it.  On the first four (exact) points of
the $r=0.5$ row the error is flat, and we announced that the decline was a near-threshold
phenomenon.  It is not: those four points lie before the summit.  Reading a pre-summit segment
as a plateau is precisely the error being corrected in this section, and four exact values
persuaded us of it more readily than eight estimated ones would have.  Exactness of values is
not length of series.

\subsection*{Why the growth form cannot be chosen}
If $E=\Theta(n)$ the constant of Section~3 is untouched, since the terms balanced there are
$\Theta(n\ln n)$.  If $E=\Theta(n\ln n)$ it moves.  The data do not decide, and cannot.

Fitting the full $r=0.5$ row to $E=c+an+bn\ln n$ returns $b$ at hundreds of standard deviations
from zero --- and a $\chi^2$ of $240$ on eight degrees of freedom, which rejects the model and
voids the uncertainty with it.  Fifteen further one-term variants using $\ln n$, $\sqrt n$ and
$n/\ln n$ are also rejected, the best by a factor of eighty-four.  The per-point error was measured rather than assumed, and it is not constant: eight fully
independent seeds give $0.262\%$ at $n=30$, $0.404\%$ at $n=35$ and $0.979\%$ at $n=40$.  At
$n=40$ the batch spread \emph{is} the error; at $n=30$ it overstates it fourfold.  A single
measured value applied to a whole series --- which both of us did, in opposite directions ---
inflates or deflates every $\chi^2$ that follows.  All figures below use the per-point values.

The rejection is confined to small $n$.  Taken on tails, three-parameter forms are admissible
from $n\ge15$ onward: $\chi^2=2.5/3$, $2.2/2$, $0.05/1$ for $c+an+bn\ln n$ at $n\ge15,20,25$.
The asymptotic regime appears to begin near $n=15$, and below it we were fitting asymptotic forms
to pre-asymptotic data.

On that tail the fit is emphatic: $b=-0.0909\pm0.0011$, $-0.0898\pm0.0024$, $-0.0824\pm0.0056$
across three nested tails, while the purely linear $c+an$ is rejected at $\chi^2=6573/4$,
$1379/3$, $216/2$.  The significance of $b$ is $81$, $37$ and $15$ standard deviations on the
three tails --- itself a warning, since a robust quantity does not lose three quarters of its
significance when four points are dropped.  Read alone this says $E=\Theta(n\ln n)$ and the constant
moves.  It does not say that.  On the same tail,
\begin{center}
\begin{tabular}{lll}
\toprule
form ($n\ge20$) & $\chi^2$ (2 d.o.f.) & implication\\
\midrule
$c+an+b\,n\ln n$ & $2.19$ & constant moves\\
$c+an+d\sqrt n$ & $0.18$ & constant intact\\
$c+an+e\,n/\ln n$ & $1.36$ & constant intact\\
$c+an+f\ln n$ & $0.44$ & constant intact\\
\bottomrule
\end{tabular}
\end{center}
All four are admissible; three imply the opposite conclusion and fit better.  An independent
implementation, run without sight of these figures, obtains reduced $\chi^2$ of $8.1$, $7.0$,
$7.4$, $7.5$ for the same four forms --- different in level, identical in verdict: mutually
indistinguishable, oppositely answering.  The eighty-six
standard deviations are a quantity computed \emph{inside} a model, and such a quantity does not
transfer to a choice \emph{between} models.  The test to apply is not whether a coefficient
differs from zero, but whether a second admissible model exists in which it is absent.  Here
several do.

Nor is this curable by computing further.  The two candidate bases are collinear to $0.9983$ over
$n\le40$, and the collinearity does not improve with range: $0.9986$ to $n=400$, $0.9993$ to
$n=10^4$.  What separates them is precision, not length, and the precision is already $0.26\%$.

We therefore withdraw the support that the earlier Section~6 offered and do not replace it.  The
derivation of Section~3 stands, but it establishes what the first moment gives, not what the
problem gives, and the distance between those two is what this section failed to measure.

\section{Where the reach ends, and why}
The limit is not an artefact of the estimator.  A descent guided by importance sampling reaches
only the depths a greedy search reaches, and the quality of greedy search degrades with $n$: we
measured the fraction of random restarts attaining the known maximum as $8/8,\ 3/8,\ 2/8,\ 1/8$
at $n=5,\dots,8$, a geometric decay.  At $n=20$ greedy attains about $0.925$ of the maximum, so
the ratio $0.9$ is within reach; by $n=25$ it falls below $0.9$ and the required depth is outside
what such a descent visits at all.  Every top-down greedy method shares this bound.  To go past it
one needs a descent guided by something that knows the target --- which is to say, the upper bound
that the problem does not have.

\section{The objection this answers}
Kaplan's objection to the Guy--Kelly conjecture, as summarised in \cite{Kaplan}, is not that the
constant is wrong but that the derivation is unsafe: it ``suggests that certain random events would
be largely independent of each other when they're really not'', and ``if we look more closely at
how non-independent they actually are we might not reach the same conclusion''.  No alternative
constant is proposed.

The objection is exactly right about the derivation, and Sections 3 and 4 quantify how badly:
the independence assumption is off by seven orders of magnitude at $n=20$, and the crude form
survives only because two of its errors cancel.  We had believed Section~6 supplied the
measurement the objection asks for.  It does not, and we now believe no measurement of this kind
can: the two possible answers correspond to bases that stay collinear at every range, and both
admit an acceptable fit to the data we can obtain.  To ``might we not reach the same conclusion''
the honest answer is that we cannot tell by counting, and that we have measured how far from
telling we are.

\section{On three dimensions}
The same test was carried out in the cube (no four coplanar) on complete enumerations at
$n=3,4,5$ by the first solver.  There the increment of the error \emph{grows}: $+2.22$, then
$+5.53$.  By the criterion above that is a wrong exponent.  The conclusion of Section 5 is
therefore two-dimensional and must not be carried into the cube.  This is consistent with the
mechanism: in the plane the hard ceiling $2n$ binds and anchors the estimate, while in the cube the
corresponding bound $3n$ stands far from the achievable values and does not.

\begin{thebibliography}{9}
\bibitem{GuyKelly} R.~K.~Guy and P.~A.~Kelly, The no-three-in-line problem, \emph{Canad. Math. Bull.} 11 (1968) 527--531.
\bibitem{Voutier} P.~M.~Voutier, On the Guy--Kelly conjecture for the no-three-in-line problem, arXiv:2603.00215 (2026).
\bibitem{Prellberg} T.~Prellberg, Constraint satisfaction programming for the no-three-in-line problem, \emph{J. Combin. Theory} 225 (2027) 106244.
\bibitem{Survey} The general position problem: a survey, arXiv:2501.19385 (2025).
\bibitem{Kaplan} D.~Eppstein, Random no-three-in-line sets (summary of a talk by H.~Kaplan), \url{https://11011110.github.io/blog/2018/11/10/random-no-three.html}.
\bibitem{Flammenkamp} A.~Flammenkamp, The no-three-in-line problem, \url{https://wwwhomes.uni-bielefeld.de/achim/no3in/readme.html}.
\end{thebibliography}

\section*{AI/tool-use disclosure}
The computations, the programs and the text were produced by autonomous AI agents (Anthropic Claude)
under the direction of the author of record, who posed the question, chose what to compute, decided
what to claim and takes responsibility for the content.  Two agents worked independently, one in
each dimension, and each checked the other's conclusions rather than its code.  Programs, journals
and the full record --- including the measurements that refuted our own earlier claims --- are at
\url{https://github.com/iwasborninbali/saturation}.
\end{document}
